\documentclass[11pt]{article}
\usepackage[margin=1in]{geometry}
\usepackage{amsmath,amssymb,amsthm,mathtools}
\usepackage{hyperref}

\newtheorem{theorem}{Theorem}[section]
\newtheorem{lemma}[theorem]{Lemma}
\newtheorem{proposition}[theorem]{Proposition}
\newtheorem{corollary}[theorem]{Corollary}
\theoremstyle{definition}
\newtheorem{convention}[theorem]{Convention}
\newtheorem{definition}[theorem]{Definition}
\theoremstyle{remark}
\newtheorem{remark}[theorem]{Remark}

\newcommand{\Z}{\mathbb{Z}}
\newcommand{\Q}{\mathbb{Q}}
\newcommand{\F}{\mathbb{F}}
\newcommand{\hgt}{\operatorname{ht}}
\newcommand{\Cset}{\mathcal{C}}
\newcommand{\one}{\mathbb{1}}
\newcommand{\cube}{\{0,1\}}

\title{Ribbon operators that add a single shape and still commute\\
{\large the structure of the vanishing locus (gap (H2) of Q150)}}
\author{Clio}
\date{16 September 2026 --- prove cycle 2}

\begin{document}
\maketitle

\begin{abstract}
The morning session \cite{Q150c1} refuted the nonvanishing lemma of \cite{Q149} and replaced
the classification of commuting pairs $[R_e^{W},R_f^{\bar W}]=0$ by a commuting-words theorem,
reducing everything to a single rank-$d$ weight $\tau$, $d=\gcd(e,f)$, subject to the two-bead
condition \textnormal{(2B$^*_d$)}.  It left open (H2): \emph{which} $\tau$ with zeros satisfy it.

This paper settles the extreme case completely and the general case halfway.

\textbf{(A)}  For a weight supported on a \emph{single} word $w\in\cube^{d-1}$ --- an operator
that adds exactly one ribbon shape and nothing else --- \textnormal{(2B$^*_d$)} holds if and
only if
\[
   w_{p+1}\;\ne\;w_{n-p}\qquad\text{for every border length }p\text{ of }w,\qquad n=d-1 .
\]
The border $p=0$ alone reads: \emph{exactly one of the two end rows of the ribbon has length
$1$}.  Consequently, for every $d\ge3$, every such $w$ and every $k\ge2$, the operator
$R_d^{\one_w}$ --- which adds the single ribbon shape $\alpha(w)\models d$ --- commutes with an
explicit $R_{kd}^{\bar W}$ supported on $2^{k-1}$ shapes.  There are $2,4,8,12,28,48,104$ such
shapes for $d=3,\dots,9$.  This strictly sharpens \cite[Thm.~6.6]{Q150c1}, whose weights are
supported on $2^{d-3}$ shapes.

\textbf{(B)}  For the strong condition \textnormal{(2B$^{**}_d$)} the single-shape solutions are
exactly the \emph{antipalindromes} $w_i\ne w_{n+1-i}$; in particular there are none for even $d$.

\textbf{(C)}  In general: every support of a \textnormal{(2B$^*_d$)}-solution is empty, everything,
or contained in one of the cylinders $\Cset_{j,b}$ of \cite[Conv.~6.4]{Q150c1} --- indeed the
cylinders are exactly the maximal proper supports --- and for the minimal such $j$ the support is
invariant under flipping every coordinate in $\{1,\dots,j-1\}\cup\{d-j+1,\dots,d-1\}$.  This
proves the containment half of (H2), which \cite{Q150c1} had only observed computationally for
$d\le5$.  Two things follow.  It explains the threshold $\gcd\ge3$: $d\ge3$ is precisely when
$[1,d-1]$ contains an index $j\ne d/2$.  And it turns the refuted lemma into a genuine
\emph{dichotomy}: a commuting $W$ either vanishes nowhere, or vanishes at more than three
quarters of the ribbon shapes of rank $d$ --- at all but a fraction $4^{-e/d}$ of those of rank
$e$.  There is nothing in between.

Independently of all this, the $(3,6)$ counterexample of \cite[Thm.~6.7]{Q150c1} is re-verified
here from the operator definition by a from-scratch implementation, and one omission in its
statement is corrected: the value $W_{(1,2)}$, which that theorem does not specify, must be $0$.
\end{abstract}

\tableofcontents

\section{What this paper does, and what it does not}

Let $\lambda$ be a partition, $M(\lambda)$ its Maya set, and for a \emph{shape weight}
$W:\cube^{e-1}\to K$ let $R_e^{W}$ be the operator of \cite[Conv.~1.2]{Q149},
\[
   R_e^{W}\,|M\rangle \;=\; \sum_{\substack{b\in M,\ b+e\notin M}}
       W\bigl(u_M(b,e)\bigr)\;\bigl|(M\setminus\{b\})\cup\{b+e\}\bigr\rangle ,
   \qquad u_M(b,e)_i=[\,b+i\in M\,].
\]
By \cite[Lemma 1.3]{Q149} the word $u$ names the shape of the added ribbon: $u\mapsto\alpha(u)$
is a bijection onto compositions of $e$, the parts being the row lengths read top row first.

\cite{Q150c1} proved that $[R_e^{W},R_f^{\bar W}]=0$ with $W\not\equiv0\not\equiv\bar W$ forces,
in the sign-free coordinates $\sigma(u)=(-1)^{|u|}W(u)$,
\[
   \sigma=c\,\tau^{\otimes(e/d)},\qquad \bar\sigma=\bar c\,\tau^{\otimes(f/d)},\qquad d=\gcd(e,f),
\]
for a single rank-$d$ weight $\tau$, and that the whole remaining content is the two-bead
condition \textnormal{(2B$^*_e$)} on $\sigma$.  Question (H2) of \cite[\S8]{Q150c1} is: what are
the rank-$d$ solutions with zeros?

The present paper answers this for the extreme case --- the smallest possible support --- and
proves the structural half of the general case.  What it does \emph{not} do is give a closed
description of the admissible supports \emph{inside} a cylinder $\Cset_{j,b}$; \S\ref{sec:gaps}
states precisely what remains, with the reduction that Theorem~\ref{thm:struct} performs.

A remark on why the single-shape case is worth isolating.  The refuted lemma
\cite[Lemma 8.2]{Q149} said that a commuting partner forces \emph{every} weight value to be
nonzero.  The sharpest possible refutation is not ``some value vanishes'' but ``all but one
value vanishes'': an operator that adds a single ribbon shape.  Theorem~\ref{thm:singleton}
says exactly which single shapes do this, and the answer is a condition on the word's borders
--- again a free-monoid statement, as was the $\gcd$ in \cite[Thm.~1.2]{Q150c1}.

\section{Setting}\label{sec:setup}

We use \cite[\S2]{Q150c1} verbatim and do not restate its proofs.  Throughout $K$ is an
integral domain, $d\ge1$, and $n:=d-1$.

\begin{convention}[{\cite[Conv.~2.1]{Q150c1}}]\label{conv:2B}
A shape weight $\sigma$ of rank $m\ge2$ satisfies \textbf{(2B$^*_m$)} if for every
$\delta$ with $1\le\delta\le m-1$ and all $\pi,\varsigma\in\cube^{\delta-1}$,
$\pi'\in\cube^{m-\delta-1}$,
\begin{equation}\label{eq:2Bstar}
  \sigma(\pi\,0\,\pi')\,\sigma(\pi'\,0\,\varsigma)
  \;=\;\sigma(\pi\,1\,\pi')\,\sigma(\pi'\,1\,\varsigma),
\end{equation}
the marked letter sitting at position $\delta$ in the first factor and at position $m-\delta$
in the second.  It satisfies the stronger \textbf{(2B$^{**}_m$)} if the same holds with the
second factor's blocks decoupled: for all $\alpha,\beta'\in\cube^{\delta-1}$ and
$\alpha',\beta\in\cube^{m-\delta-1}$,
\begin{equation}\label{eq:2Bss}
  \sigma(\alpha\,0\,\alpha')\,\sigma(\beta\,0\,\beta')
  \;=\;\sigma(\alpha\,1\,\alpha')\,\sigma(\beta\,1\,\beta') .
\end{equation}
For $m=1$ both are vacuous.
\end{convention}

Two facts from \cite{Q150c1} are used as stated.

\begin{itemize}
\item \cite[Thm.~1.3]{Q150c1} (the classification): for $1\le e<f$, $d=\gcd(e,f)$, over a field
$F$, $[R_e^{W},R_f^{\bar W}]=0$ with both weights nonzero \emph{iff} there are a nonzero
rank-$d$ weight $\tau$ and $c,\bar c\in F^\times$ with $\sigma=c\,\tau^{\otimes(e/d)}$,
$\bar\sigma=\bar c\,\tau^{\otimes(f/d)}$ and $\sigma$ satisfies (2B$^*_e$).
\item \cite[Lemma 5.1]{Q150c1} (descent), part (a) in the case $k=1$: if $e=d$ then
$\sigma=\tau$, and (2B$^*_d$) for $\tau$ is exactly (2B$^*_e$) for $\sigma$.
\end{itemize}

\begin{definition}[support condition]\label{def:S}
For $Z\subseteq\cube^{n}$ say that $Z$ is \textbf{admissible} (for rank $d=n+1$) if for every
$\delta\in[1,d-1]$ and all $\pi,\varsigma\in\cube^{\delta-1}$, $\pi'\in\cube^{d-\delta-1}$,
\begin{equation}\label{eq:S}
  \bigl[\pi0\pi'\in Z\ \wedge\ \pi'0\varsigma\in Z\bigr]
  \iff
  \bigl[\pi1\pi'\in Z\ \wedge\ \pi'1\varsigma\in Z\bigr].
\end{equation}
\end{definition}

\begin{proposition}\label{prop:supp}
Let $\sigma$ be a shape weight of rank $d$ over an integral domain and
$Z=\{u:\sigma(u)\ne0\}$.  If $\sigma$ satisfies \textnormal{(2B$^*_d$)} then $Z$ is admissible.
Conversely, if $Z$ is admissible then the indicator $\one_Z$ satisfies
\textnormal{(2B$^*_d$)}.  Hence the admissible subsets of $\cube^{n}$ are exactly the supports
of \textnormal{(2B$^*_d$)}-solutions.
\end{proposition}

\begin{proof}
In a domain a product of two elements is $0$ iff a factor is $0$, so
$\sigma(\pi0\pi')\sigma(\pi'0\varsigma)\ne0$ iff both arguments lie in $Z$; equality of the two
sides of \eqref{eq:2Bstar} therefore forces \eqref{eq:S}.  Conversely for $\sigma=\one_Z$ each
side of \eqref{eq:2Bstar} is $1$ if both its arguments lie in $Z$ and $0$ otherwise, and
\eqref{eq:S} says these two alternatives agree.
\end{proof}

Proposition~\ref{prop:supp} is what makes the support question self-contained: no separate
existence argument for the values is needed, and every admissible $Z$ really occurs.

\section{Independent verification of the $(3,6)$ counterexample, and one correction}
\label{sec:verify36}

Before building on \cite{Q150c1} I re-derived its load-bearing object from
\cite[Conv.~1.2]{Q149} alone, in a fresh implementation that shares no code with the Q149 or
Q150 engines (\texttt{code/2026-09-16-c2-Q150-H2/independent\_maya.py}).  Maya sets are stored
as $M=\{n<-L\}\cup S$ and the operator is applied literally as written.

\paragraph{Calibration.}  The dictionary \cite[Lemma 1.3]{Q149} was re-checked against Young
diagrams: for every partition $\lambda$ with $|\lambda|\le8$, every $e\le5$ and every legal bead
move, the row lengths of $\mu/\lambda$ were compared with $\alpha(u_M(b,e))$.  All $1375$ moves
agree (the same population size as \cite[Check V1]{Q149}, reached independently).

\paragraph{The counterexample.}  With $W$ and $\bar W$ of \cite[Thm.~6.7]{Q150c1},
$[R_3^{W},R_6^{\bar W}]$ annihilates every one of the $195$ partitions with $|\lambda|\le11$,
and both operators act nontrivially ($R_3^{W}$ on $21$ partitions of size $\le6$, $R_6^{\bar W}$
on $17$ of size $\le8$).  Seven perturbations were run as negative controls; all fire.  A
Murnaghan--Nakayama positive control ($\sigma\equiv1$) is silent, as it must be.  Details in
\S\ref{sec:ver}, V1--V2.

\paragraph{A correction to \cite[Thm.~6.7]{Q150c1}.}  That theorem specifies
$W_{(2,1)}=-1$ and $W_{(3)}=W_{(1,1,1)}=0$, and says nothing about the fourth composition of
$3$.  The value is not free: control NC1 sets $W_{(1,2)}=1$ and the commutator becomes nonzero
on $35$ of the $97$ partitions with $|\lambda|\le9$.  The intended weight is of course
$\sigma=\one_{(1,0)}$, which gives $W_{(1,2)}=0$ --- as the abstract of \cite{Q150c1} says in
words (``adds \emph{only} the L-tromino $(2,1)$'').  The statement should read
\[
  W_{(2,1)}=-1,\qquad W_{(3)}=W_{(1,2)}=W_{(1,1,1)}=0 .
\]
Nothing else in \cite{Q150c1} is affected: its proof constructs $\tau$ as the indicator of
$\Cset_{1,1}=\{(1,0)\}$ and therefore has the right object throughout.  Only the displayed
statement is incomplete.

\section{Single-shape weights}\label{sec:singleton}

\begin{definition}\label{def:border}
Let $w\in\cube^{n}$.  An integer $p$ with $0\le p\le n-1$ is a \textbf{border length} of $w$ if
$w_{[1,p]}=w_{[n-p+1,n]}$.  (Thus $p=0$ is always a border length.)  Say $w$ is
\textbf{non-overlapping} if
\begin{equation}\label{eq:NB}\tag{NB}
   w_{p+1}\;\ne\;w_{n-p}\qquad\text{for every border length }p\text{ of }w .
\end{equation}
\end{definition}

In words: for each way in which a prefix of $w$ coincides with the suffix of the same length,
the letter immediately \emph{after} the prefix occurrence and the letter immediately
\emph{before} the suffix occurrence must differ.

\begin{theorem}[single-shape weights, weak condition]\label{thm:singleton}
Let $d\ge2$, $n=d-1$, let $K$ be an integral domain, $c\in K\setminus\{0\}$ and
$w\in\cube^{n}$.  Then $\tau=c\,\one_{w}$ satisfies \textnormal{(2B$^*_d$)} if and only if $w$
is non-overlapping.
\end{theorem}

\begin{proof}
Fix $\delta\in[1,d-1]$ and $\gamma\in\cube$, and consider the side
$\tau(\pi\gamma\pi')\,\tau(\pi'\gamma\varsigma)$ of \eqref{eq:2Bstar}, where
$|\pi|=|\varsigma|=\delta-1$ and $|\pi'|=d-\delta-1=n-\delta$.  Since $\tau$ is supported on the
single word $w$, this side is nonzero iff both arguments equal $w$, i.e.\ iff
\[
   \pi=w_{[1,\delta-1]},\quad \gamma=w_\delta,\quad \pi'=w_{[\delta+1,n]}
   \qquad\text{and}\qquad
   \pi'=w_{[1,n-\delta]},\quad \gamma=w_{n-\delta+1},\quad \varsigma=w_{[n-\delta+2,n]},
\]
the two prescriptions of $\pi'$ having the same length $n-\delta$.  So, for suitable
$\pi,\pi',\varsigma$, the side is nonzero exactly when
\begin{equation}\label{eq:twoconds}
  w_{[\delta+1,n]}=w_{[1,n-\delta]}
  \qquad\text{and}\qquad
  \gamma=w_\delta=w_{n-\delta+1},
\end{equation}
and then its value is $c^2$.

Substitute $p:=n-\delta$, so that $p$ runs over $[0,n-1]$ as $\delta$ runs over $[1,n]=[1,d-1]$.
The first condition in \eqref{eq:twoconds} reads $w_{[n-p+1,n]}=w_{[1,p]}$, i.e.\ $p$ is a
border length of $w$; the second reads $\gamma=w_{n-p}=w_{p+1}$.

Suppose $w$ is not non-overlapping, and let $p$ be a border length with $w_{p+1}=w_{n-p}=:\gamma_0$.
Take $\delta=n-p$ and the arguments prescribed above.  Then the $\gamma_0$-side of
\eqref{eq:2Bstar} equals $c^2$, while the $(1-\gamma_0)$-side is $0$ because
$\gamma=1-\gamma_0$ violates the second condition of \eqref{eq:twoconds}.  As $K$ is a domain
and $c\ne0$ we have $c^2\ne0$, so \eqref{eq:2Bstar} fails.  (This is the only appeal to $K$
being a domain in the proof, and its witness is the single element $c$.)

Conversely, suppose $w$ is non-overlapping.  Then for every $\delta$ and every $\gamma$ the
conditions \eqref{eq:twoconds} fail, so both sides of \eqref{eq:2Bstar} vanish identically and
$\tau$ satisfies \textnormal{(2B$^*_d$)}.
\end{proof}

\begin{remark}[reading \eqref{eq:NB} on the ribbon]\label{rem:endrows}
The border $p=0$ contributes the condition $w_1\ne w_n$.  By \cite[Lemma 1.3]{Q149}, $w_1=1$ iff
the \emph{last} part of $\alpha(w)$ is $1$ and $w_{n}=1$ iff the \emph{first} part is $1$.  So
the $p=0$ clause says: \emph{exactly one of the two end rows of the ribbon has length $1$}.
For $d\le5$ this clause is the whole of \eqref{eq:NB}; the first $w$ satisfying it but failing
\eqref{eq:NB} occurs at $d=6$, namely $w=01001$ (border $p=2$, with $w_3=w_3$), which has
period $3=d/2$.  Verified for $d\le7$ (Check V4).
\end{remark}

\begin{remark}[$d\le2$]
For $d=2$, $n=1$: the only border length is $p=0$ and the condition is $w_1\ne w_1$, false.
So there are no single-shape solutions at $d\le2$, consistent with
\cite[Thm.~6.1]{Q150c1} (nonvanishing for $\gcd\le2$).
\end{remark}

\begin{theorem}[single-shape weights, strong condition]\label{thm:singleton2}
In the situation of Theorem~\ref{thm:singleton}, $\tau=c\,\one_w$ satisfies
\textnormal{(2B$^{**}_d$)} if and only if $w$ is an \textbf{antipalindrome},
$w_i\ne w_{n+1-i}$ for all $i$.  In particular there is no such $w$ when $n$ is odd, i.e.\ when
$d$ is even; and for $d$ odd there are exactly $2^{(d-1)/2}$ of them.
\end{theorem}

\begin{proof}
Fix $\delta$ and $\gamma$ and consider $\tau(\alpha\gamma\alpha')\tau(\beta\gamma\beta')$ with
$|\alpha|=|\beta'|=\delta-1$, $|\alpha'|=|\beta|=n-\delta$.  It is nonzero iff
$\alpha\gamma\alpha'=w$ and $\beta\gamma\beta'=w$.  In the first word $\gamma$ sits at position
$|\alpha|+1=\delta$, in the second at position $|\beta|+1=n-\delta+1$; the remaining blocks are
then determined and unconstrained.  So the $\gamma$-side is nonzero iff
$\gamma=w_\delta=w_{n+1-\delta}$, with value $c^2$.  Hence the two sides of \eqref{eq:2Bss}
agree iff both vanish, iff $w_\delta\ne w_{n+1-\delta}$.  Ranging over $\delta\in[1,n]$ gives
the antipalindrome condition.  If $n$ is odd, $\delta=(n+1)/2$ demands $w_\delta\ne w_\delta$;
if $n$ is even the condition determines $w_{[n/2+1,n]}$ freely from $w_{[1,n/2]}$.
\end{proof}

\begin{remark}
An antipalindrome is non-overlapping: if $p$ is a border length then, comparing
$w_{p+1}$ with $w_{n-p}=w_{n+1-(p+1)}$, the antipalindrome condition at $i=p+1$ gives
$w_{p+1}\ne w_{n-p}$ directly.  So Theorem~\ref{thm:singleton2} describes a subfamily of
Theorem~\ref{thm:singleton}, as it must, and the inclusion is strict for $d\ge4$
(e.g.\ $d=4$: no antipalindromes, four non-overlapping words).  This is one more instance of
\cite[Rem.~5.2]{Q150c1}: (2B$^{**}$) is strictly stronger than (2B$^*$) from $d\ge4$ on.
\end{remark}

\section{Operators that add a single ribbon shape and commute}\label{sec:counter}

\begin{theorem}\label{thm:main}
Let $d\ge3$, let $w\in\cube^{d-1}$ be non-overlapping, let $\alpha=\alpha(w)\models d$ be the
corresponding ribbon shape, let $K$ be any integral domain and let $k\ge2$.  Put $e=d$,
$f=kd$ and
\[
   W \;=\; \one_{w}\ \ (\text{rank }d),\qquad
   \bar W(y)\;=\;(-1)^{\hgt(y)}\prod_{i=1}^{k}
      \bigl[\,y_{[(i-1)d+1,\;id-1]}=w\,\bigr]\ \ (\text{rank } kd).
\]
Then:
\begin{enumerate}
\item[\textup{(i)}] $R_d^{W}s_\lambda=\sum_\mu s_\mu$, the sum over those $d$-ribbons
  $\mu/\lambda$ of shape exactly $\alpha$; the operator adds one ribbon shape and nothing else.
\item[\textup{(ii)}] $\bar W$ is supported on exactly $2^{k-1}$ of the $2^{kd-1}$ shapes of rank
  $kd$, namely on the words $w\,c_1\,w\,c_2\cdots c_{k-1}\,w$ with $c_i\in\cube$ free.
\item[\textup{(iii)}] $[R_d^{W},R_{kd}^{\bar W}]=0$.
\end{enumerate}
Since $\gcd(d,kd)=d\ge3$, each such pair is a counterexample to \cite[Lemma 8.2]{Q149}, and
$W$ has the smallest support any counterexample can have.
\end{theorem}

\begin{proof}
(i) is the definition of $R_d^{W}$ together with \cite[Lemma 1.3]{Q149}, and (ii) is the
definition of $\tau^{\otimes k}$ (\cite[Conv.~1.1]{Q150c1}): $\sigma(u)=\prod_i
\tau(u_{[(i-1)d+1,id-1]})$, the $k-1$ letters in positions $d,2d,\dots,(k-1)d$ being ignored.

For (iii) put $\tau:=\one_w$, so that $\sigma=\tau$ (as $e/d=1$) and $\bar\sigma=\tau^{\otimes k}$,
and note $W=(-1)^{\hgt(w)}\sigma$ pointwise; rescaling $W$ by the unit $(-1)^{\hgt(w)}$ does not
affect whether the commutator vanishes.  By Theorem~\ref{thm:singleton}, $\tau$ satisfies
(2B$^*_d$), which for $e=d$ is (2B$^*_e$) for $\sigma$.  Now apply
\cite[Thm.~1.3]{Q150c1}, direction $(\Leftarrow)$, with $c=\bar c=1$: the hypotheses are
$\sigma=\tau^{\otimes(e/d)}$, $\bar\sigma=\tau^{\otimes(f/d)}$ and (2B$^*_e$) for $\sigma$, all
of which hold.  Hence the commutator vanishes.

Finally, a weight with empty support is $\equiv0$, which \cite[Lemma 8.2]{Q149} excludes by
hypothesis; so a single point is the smallest possible support.
\end{proof}

\begin{corollary}[the smallest new case]\label{cor:48}
Let $K$ be any integral domain.  Let $W$ be the rank-$4$ shape weight with
\[
  W_{(1,3)}=-1,\qquad W_\beta=0\ \text{ for the other seven compositions }\beta\models4,
\]
and let $\bar W$ be the rank-$8$ weight with
\[
  \bar W_{(1,4,3)}=+1,\qquad \bar W_{(1,3,1,3)}=-1,\qquad \bar W_\beta=0\ \text{ for the other }
  126\text{ compositions }\beta\models8 .
\]
Then $[R_4^{W},R_8^{\bar W}]=0$ and both operators are nonzero.
\end{corollary}

\begin{proof}
Theorem~\ref{thm:main} with $d=4$, $k=2$, $w=(0,0,1)$: the border lengths of $001$ are $p=0$
only, and $w_1=0\ne1=w_3$, so $w$ is non-overlapping; $\alpha(001)=(1,3)$.  The support of
$\tau^{\otimes2}$ is $\{001\,c\,001:c\in\cube\}=\{0010001,0011001\}$, with
$\alpha(0010001)=(1,4,3)$ and $\alpha(0011001)=(1,3,1,3)$, of heights $2$ and $3$, so
$(-1)^{\hgt}$ is $+1$ and $-1$ respectively; and $\hgt(001)=1$ gives $W_{(1,3)}=-1$.  Verified directly in \S\ref{sec:ver}, V5.
\end{proof}

\begin{remark}
The list of admissible single shapes begins:
\[
\begin{array}{ll}
 d=3: & (1,2),\,(2,1)\\
 d=4: & (1,3),\,(3,1),\,(1,1,2),\,(2,1,1)\\
 d=5: & (1,4),\,(4,1),\,(1,1,3),\,(3,1,1),\,(1,2,2),\,(2,2,1),\,(1,1,1,2),\,(2,1,1,1)
\end{array}
\]
with $2,4,8,12,28,48,104$ shapes for $d=3,\dots,9$.  The list is stable under ribbon
transposition, which by \cite[Lemma 1.6]{Q149} is reversal-and-complementation on words: both
operations preserve border lengths and swap the two sides of \eqref{eq:NB}.  At $d=3$ it
recovers \cite[Thm.~6.7]{Q150c1} ($w=(1,0)$, shape $(2,1)$) and its transpose.
\end{remark}

\section{The structure of an admissible support}\label{sec:struct}

We now drop the single-shape hypothesis.  Recall from \cite[Conv.~6.4]{Q150c1}, for
$1\le j\le d-1$ with $2j\ne d$ and $b\in\cube$,
\[
   \Cset_{j,b}\;=\;\bigl\{\,w\in\cube^{n}\ :\ w_j=b,\ w_{d-j}=1-b\,\bigr\},
   \qquad |\Cset_{j,b}|=2^{\,d-3}.
\]
Note $\Cset_{j,b}=\Cset_{d-j,1-b}$, so there are $2\lfloor(d-1)/2\rfloor$ distinct cylinders.

\begin{theorem}\label{thm:struct}
Let $d\ge2$ and let $Z\subseteq\cube^{n}$ be admissible (Definition~\ref{def:S}).  Then either
$Z=\emptyset$, or $Z=\cube^{n}$, or there are $j\in[1,d-1]$ with $2j\ne d$ and $b\in\cube$ with
$Z\subseteq\Cset_{j,b}$.  Moreover, in the third case let $j_0$ be minimal such; then $Z$ is
invariant under flipping any single coordinate in
\[
   \{1,\dots,j_0-1\}\ \cup\ \{d-j_0+1,\dots,d-1\}.
\]
\end{theorem}

\begin{proof}
For $j\in[1,d-1]$ put $\delta:=d-j$ and, for $\pi'\in\cube^{\,j-1}$ and $b\in\cube$,
\[
  A_b(\pi')=\{\pi\in\cube^{\,d-j-1}:\ \pi\,b\,\pi'\in Z\},\qquad
  B_b(\pi')=\{\varsigma\in\cube^{\,d-j-1}:\ \pi'\,b\,\varsigma\in Z\} .
\]
The lengths match those in Definition~\ref{def:S}: $|\pi|=|\varsigma|=\delta-1=d-j-1$ and
$|\pi'|=d-\delta-1=j-1$.  In $\pi b\pi'$ the letter $b$ sits at position $\delta=d-j$; in
$\pi'b\varsigma$ it sits at position $d-\delta=j$.  Condition \eqref{eq:S} at $\delta$ says
precisely
\begin{equation}\label{eq:prodeq}
  A_0(\pi')\times B_0(\pi')\;=\;A_1(\pi')\times B_1(\pi')\qquad\text{for every }\pi' .
\end{equation}

Let $F_0=\emptyset$ and $F_j=F_{j-1}\cup\{j,\ d-j\}$, so $F_j=\{1,\dots,j\}\cup\{d-j,\dots,d-1\}$.
Write $P(j)$ for the statement ``$Z$ is invariant under flipping any single coordinate in
$F_j$''; $P(0)$ is vacuously true.

\emph{Inductive step.}  Fix $j\ge1$ and assume $P(j-1)$, i.e.\ invariance at every coordinate in
$\{1,\dots,j-1\}\cup\{d-j+1,\dots,d-1\}$.  In $\pi\,b\,\pi'$ the block $\pi'$ occupies positions
$d-j+1,\dots,d-1$, and in $\pi'\,b\,\varsigma$ it occupies positions $1,\dots,j-1$.  By $P(j-1)$,
membership in $Z$ is unchanged by flipping any of these positions, so $A_b(\pi')$ and
$B_b(\pi')$ do not depend on $\pi'$; write $A_b$, $B_b$.

Two subsets $A_0\times B_0$ and $A_1\times B_1$ of a product are equal iff either both are empty
or $A_0=A_1$ and $B_0=B_1$ (a nonempty product determines its factors as its projections).  So
\eqref{eq:prodeq} leaves exactly two cases.

\emph{Case 1: $A_0=A_1$ and $B_0=B_1$.}  $A_0=A_1$ says $\pi0\pi'\in Z\iff\pi1\pi'\in Z$ for all
$\pi,\pi'$, i.e.\ $Z$ is invariant under flipping coordinate $\delta=d-j$; $B_0=B_1$ says the
same for coordinate $j$.  Together with $P(j-1)$ this is $P(j)$.

\emph{Case 2: both products empty.}  Then ($A_0=\emptyset$ or $B_0=\emptyset$) and
($A_1=\emptyset$ or $B_1=\emptyset$).  Now $A_b=\emptyset$ says no word of $Z$ has $b$ in
position $d-j$, and $B_b=\emptyset$ says no word of $Z$ has $b$ in position $j$.  Hence:
\begin{itemize}
\item $A_0=A_1=\emptyset$, or $B_0=B_1=\emptyset$: every word of $Z$ would have to avoid both
  values at one position, so $Z=\emptyset$.
\item $A_0=\emptyset$ and $B_1=\emptyset$: every $w\in Z$ has $w_{d-j}=1$ and $w_j=0$, i.e.\
  $Z\subseteq\Cset_{j,0}$.
\item $A_1=\emptyset$ and $B_0=\emptyset$: symmetrically $Z\subseteq\Cset_{j,1}$.
\end{itemize}
If $2j=d$ the last two read $Z\subseteq\{w:w_j=0,\ w_j=1\}=\emptyset$, consistent with
$\Cset_{j,b}$ being undefined (empty) there.

\emph{Termination.}  Let $J=\lceil (d-1)/2\rceil$; then $J\le d-1$ for $d\ge2$, and
$F_J=\{1,\dots,J\}\cup\{d-J,\dots,d-1\}=\{1,\dots,d-1\}$ because the gap $\{J+1,\dots,d-J-1\}$
is empty when $2J\ge d-1$.  Run the inductive step for $j=1,2,\dots,J$.  If Case 2 never occurs
then $P(J)$ holds, so $Z$ is invariant under flipping every coordinate, hence
$Z\in\{\emptyset,\cube^{n}\}$.  Otherwise let $j_1$ be the first index at which Case 2 occurs;
then $Z=\emptyset$ or $Z\subseteq\Cset_{j_1,b}$ for some $b$.

\emph{Minimality.}  Suppose $Z\ne\emptyset$ and $Z\subseteq\Cset_{j_1,b}$.  For $j<j_1$ Case 1
held, so $Z$ is invariant under flipping coordinate $j$; but $Z\subseteq\Cset_{j,b'}$ would force
$w_j=b'$ for all $w\in Z$, incompatible with that invariance unless $Z=\emptyset$.  Hence no
$j<j_1$ has $Z\subseteq\Cset_{j,b'}$, so $j_0=j_1$, and $P(j_0-1)$ is exactly the asserted
invariance.
\end{proof}

\begin{corollary}[why the threshold is $3$]\label{cor:three}
For $d\le2$ the only admissible sets are $\emptyset$ and $\cube^{n}$; a rank-$d$ weight
satisfying \textnormal{(2B$^*_d$)} and $\not\equiv0$ therefore vanishes nowhere.  For $d\ge3$
proper nonempty admissible sets exist, e.g.\ $\Cset_{1,b}$.
\end{corollary}

\begin{proof}
For $d=1$ there is nothing to prove ($n=0$).  For $d=2$ the only index is $j=1=d/2$, so
Theorem~\ref{thm:struct} leaves only $\emptyset$ and $\cube^{1}$; combined with
Proposition~\ref{prop:supp} this gives the claim.  For $d\ge3$ the index $j=1$ satisfies
$2j\ne d$, and $\Cset_{1,b}$ is admissible by \cite[Lemma 6.5]{Q150c1}.
\end{proof}

This is the conceptual content of the threshold ``$\gcd(e,f)\ge3$'' in \cite[Thm.~1.4]{Q150c1}:
$d\ge3$ is exactly the condition that the index set $[1,d-1]$ contains some $j$ with
$j\ne d-j$.  Corollary~\ref{cor:three} reproves the support half of \cite[Thm.~6.1]{Q150c1}
(nonvanishing for $\gcd\le2$) without any of its case analysis.

\begin{corollary}[the maximal proper supports are exactly the cylinders]\label{cor:maximal}
Let $d\ge3$.  The maximal proper admissible subsets of $\cube^{n}$ are precisely the
$2\lfloor(d-1)/2\rfloor$ distinct cylinders $\Cset_{j,b}$.
\end{corollary}

\begin{proof}
Each $\Cset_{j,b}$ is admissible \cite[Lemma 6.5]{Q150c1} and proper, having
$2^{d-3}<2^{d-1}$ elements.  By Theorem~\ref{thm:struct} every proper admissible set is contained
in some cylinder, so every maximal one is a cylinder.  Distinct cylinders all have the same
cardinality $2^{d-3}$, so none strictly contains another, and each is therefore maximal.
Verified for $d\le7$ (Check V3).
\end{proof}

\begin{corollary}[the quarter dichotomy]\label{cor:quarter}
Let $\tau\not\equiv0$ be a rank-$d$ shape weight over an integral domain satisfying
\textnormal{(2B$^*_d$)}.  Then \emph{either} $\tau$ vanishes nowhere, \emph{or}
\[
   \#\{u:\tau(u)\ne0\}\;\le\;2^{\,d-3}\;=\;\tfrac14\cdot2^{\,d-1},
\]
so $\tau$ vanishes at more than three quarters of all ribbon shapes of rank $d$.  There is
nothing in between.

Consequently, if $[R_e^{W},R_f^{\bar W}]=0$ with $W\not\equiv0\not\equiv\bar W$ and $d=\gcd(e,f)$,
then either $W$ is nowhere zero, or $W$ is supported on at most a fraction $4^{-e/d}$ of the
$2^{e-1}$ ribbon shapes of rank $e$ --- and likewise for $\bar W$ with $f$ in place of $e$.
\end{corollary}

\begin{proof}
The first statement is Proposition~\ref{prop:supp} and Corollary~\ref{cor:maximal}: the support
is admissible, hence empty (excluded), everything, or inside a cylinder, which has $2^{d-3}$
elements.  For the second, \cite[Thm.~1.3]{Q150c1} gives $\sigma=c\,\tau^{\otimes k}$ with
$k=e/d$, and $|\operatorname{supp}\tau^{\otimes k}|=|\operatorname{supp}\tau|^{k}\cdot2^{\,k-1}$,
the $k-1$ ignored letters being free.  If $\tau$ is nowhere zero this is
$2^{k(d-1)+k-1}=2^{e-1}$; otherwise it is at most $2^{k(d-3)+k-1}=2^{e-1-2k}$.
\end{proof}

This is the shape the statement should have taken from the start: not ``no value vanishes'' but
``either none vanishes or almost all do''.  The refuted \cite[Lemma 8.2]{Q149} asserted the first
alternative of a genuine dichotomy and lost the second.

\begin{remark}[admissibility is a Horn condition]\label{rem:horn}
Theorem~\ref{thm:closure} exhibits \eqref{eq:S} as a conjunction of definite Horn clauses
$x_u\wedge x_v\to x_{u'}$ in the $2^{n}$ propositional variables $x_w=[\,w\in Z\,]$, and
Proposition~\ref{prop:moore} is the classical fact that a class of subsets is Horn-definable iff
it is closed under intersection.  This is, I think, the right expectation to have about (H2a)
below: counting the models of a Horn theory is not the sort of thing that has a product formula,
and the counts $1,3,11,31,356$ of \S\ref{sec:gaps} do not look like one.
\end{remark}

\begin{remark}[the symmetry group]\label{rem:sym}
Admissibility is preserved by complementing all letters ($w\mapsto\bar w$) and by reversal
($w\mapsto w^{\mathrm{rev}}$).  Complementation exchanges the two sides of \eqref{eq:S};
reversal carries the relation at $\delta$ to the relation at $d-\delta$ with the roles of the
two factors exchanged, the shared block $\pi'$ going to its reverse.  The composite is ribbon
transposition (\cite[Lemma 1.6]{Q149}).  Accordingly $\Cset_{j,b}$ is carried to
$\Cset_{j,1-b}$ and $\Cset_{d-j,b}=\Cset_{j,1-b}$ respectively.
\end{remark}

\section{Generation: the admissible sets are a closure system}\label{sec:closure}

Theorem~\ref{thm:struct} says where an admissible set can live.  This section says how to
build one.  Both statements are short; together they are what I can offer in place of the closed
classification that (H2) asks for.

\begin{definition}\label{def:linked}
Let $u,v\in\cube^{n}$ and $\delta\in[1,n]$.  The pair $(u,v)$ is \textbf{$\delta$-linked} if
\[
   u_{[\delta+1,n]}=v_{[1,n-\delta]}\qquad\text{and}\qquad u_\delta=v_{n-\delta+1}.
\]
For such a pair write $\iota_\delta(u,v):=\bigl(u^{(\delta)},\,v^{(n-\delta+1)}\bigr)$, where
$x^{(i)}$ is $x$ with the letter in position $i$ flipped.
\end{definition}

Flipping position $\delta$ of $u$ does not touch $u_{[\delta+1,n]}$ and flipping position
$n-\delta+1$ of $v$ does not touch $v_{[1,n-\delta]}$, so $\iota_\delta$ carries $\delta$-linked
pairs to $\delta$-linked pairs, and it is an involution.

\begin{theorem}[generation]\label{thm:closure}
$Z\subseteq\cube^{n}$ is admissible if and only if it is closed under the rule
\begin{equation}\label{eq:rule}\tag{R}
   u,v\in Z\ \text{$\delta$-linked}\quad\Longrightarrow\quad
   u^{(\delta)}\in Z\ \text{ and }\ v^{(n-\delta+1)}\in Z .
\end{equation}
\end{theorem}

\begin{proof}
Fix $\delta$ and a triple $(\pi,\pi',\varsigma)$ as in Definition~\ref{def:S}, and put
$u_\gamma=\pi\gamma\pi'$, $v_\gamma=\pi'\gamma\varsigma$.  Both $(u_0,v_0)$ and $(u_1,v_1)$ are
$\delta$-linked, and $\iota_\delta(u_0,v_0)=(u_1,v_1)$.  Conversely a $\delta$-linked pair
$(u,v)$ comes from exactly one triple, namely $\pi=u_{[1,\delta-1]}$,
$\pi'=u_{[\delta+1,n]}=v_{[1,n-\delta]}$, $\varsigma=v_{[n-\delta+2,n]}$.  So, writing
$P(u,v)$ for ``$u\in Z$ and $v\in Z$'', condition \eqref{eq:S} at $(\delta,\pi,\pi',\varsigma)$
is exactly $P(u_0,v_0)\iff P(\iota_\delta(u_0,v_0))$, and \eqref{eq:rule} is the implication
$P(u,v)\Rightarrow P(\iota_\delta(u,v))$ for every $\delta$-linked pair.  Since $\iota_\delta$ is
an involution on $\delta$-linked pairs, the implication for all pairs is equivalent to the
biconditional for all pairs.
\end{proof}

\begin{proposition}\label{prop:moore}
The admissible subsets of $\cube^{n}$ are closed under arbitrary intersections, and
$\cube^{n}$ is admissible.  They therefore form a closure system: there is a closure operator
$X\mapsto\langle X\rangle$, the least admissible set containing $X$, and the admissible sets are
exactly the sets of the form $\langle X\rangle$.
\end{proposition}

\begin{proof}
Let $(Z_t)_{t\in T}$ be admissible and $Z=\bigcap_t Z_t$.  For any $u,v$, ``$u\in Z$ and $v\in Z$''
is the conjunction over $t$ of ``$u\in Z_t$ and $v\in Z_t$''.  So with $P_t$, $P$ as in the
previous proof, $P=\bigwedge_t P_t$ on both sides of \eqref{eq:S}, and a conjunction of
biconditionals is a biconditional between the conjunctions.  Admissibility of $\cube^{n}$ is
immediate (both sides of \eqref{eq:S} are always true).  The rest is the standard fact about
Moore families.
\end{proof}

\begin{remark}
Theorems~\ref{thm:singleton} and \ref{thm:closure} now read as one statement:
$\langle\{w\}\rangle=\{w\}$ if and only if $w$ is non-overlapping.  Indeed \eqref{eq:rule} applied
to $u=v=w$ is nontrivial exactly when $(w,w)$ is $\delta$-linked, i.e.\ when $w$ has period
$\delta$ and $w_\delta=w_{n-\delta+1}$ --- the failure condition of Theorem~\ref{thm:singleton}.
The closure of a single word is computed in Check V7; its size is
\[
\begin{array}{r|l}
 d=3 & 1\ (\times2),\quad 4\ (\times2)\\
 d=4 & 1\ (\times4),\quad 8\ (\times4)\\
 d=5 & 1\ (\times8),\quad 4\ (\times4),\quad 16\ (\times4)\\
 d=6 & 1\ (\times12),\ 2\ (\times4),\ 8\ (\times8),\ 32\ (\times8)
\end{array}
\]
--- so a single word generates itself, or a cylinder ($4$ at $d=5$, $8$ at $d=6$), or something
smaller, or everything.  Union is \emph{not} an admissible operation: at $d=5$ both
$\{0001\}$ and $\{0101\}$ are admissible and their union is not.
\end{remark}

\begin{remark}[where the ribbon structure enters]\label{rem:whereribbon}
It is worth naming the step, because the analogous statement is false in general.  A
multiplicative structure on $\cube^{n}$ does \emph{not} force its support to be a product set
without a positivity hypothesis --- this is why Hammersley--Clifford assumes $P>0$ --- so any
argument here that used only ``connectivity of the cube'' would prove something known to be
false.  The proof of Theorem~\ref{thm:struct} uses no connectivity.  What it uses is that the
two-bead relation \eqref{eq:2Bstar} couples position $\delta$ of one factor with position
$m-\delta$ of the \emph{other}, across a \emph{shared block} $\pi'$.  That is the ribbon input:
it is \cite[Lemma 4.3]{Q149}, where $\pi'$ is the overlap of the two bead windows.  The shared
block is what makes $A_b(\pi')$ and $B_b(\pi')$ live on a common index set, so that
\eqref{eq:prodeq} is an equality of \emph{product sets} --- and product-set rigidity, not
connectivity, is what forces the dichotomy.  The only way out is an asymmetry between positions
$j$ and $d-j$, which is exactly what $\Cset_{j,b}$ is.
\end{remark}

\section{Values on a given support}\label{sec:values}

Fix an admissible $Z$ and let $F$ be a field.  By Proposition~\ref{prop:supp} the solutions of
\textnormal{(2B$^*_d$)} with support exactly $Z$ form a torsor under the group
\[
   \mathrm{Hom}\bigl(\Z^{Z}/L_Z,\ F^\times\bigr),\qquad
   L_Z=\bigl\langle\, e_{u_0}+e_{v_0}-e_{u_1}-e_{v_1}\ :\
      u_0,v_0,u_1,v_1\in Z \,\bigr\rangle ,
\]
the generators running over the relations \eqref{eq:2Bstar} \emph{all four of whose arguments lie
in $Z$} (the others are $0=0$), with base point $\one_Z$.  The invariant factors of $L_Z$ were
computed for all admissible $Z$ at $d\le6$ (Check V6).  Sample:
\[
\begin{array}{llll}
 d=3 & |Z|=1:\ (F^\times)^1 & |Z|=4:\ (F^\times)^2\\
 d=4 & |Z|=1:\ (F^\times)^1 & |Z|=2:\ (F^\times)^2 & |Z|=8:\ (F^\times)^2\times\mu_2\\
 d=6 & |Z|=1:\ (F^\times)^1 & |Z|=3:\ (F^\times)^3 & |Z|=32:\ (F^\times)^3\times\mu_2
\end{array}
\]
where $\mu_2=\{x\in F^\times:x^2=1\}$.  Two remarks.  First, a single-shape support always gives
a one-dimensional torus, the free scalar $c$ of Theorem~\ref{thm:singleton} --- as it must,
there being no relations with all four arguments in a singleton unless
$\pi0\pi'=\pi1\pi'$, which is impossible.  Second, $2$-torsion does occur (already at $d=4$ on
the full support), so the solution set is not a torus over every field and in characteristic
$2$ some sign choices collapse; this matches \cite[Cor.~5.4]{Q149}, whose group has a $\mu_2$
factor.

\section{Verification}\label{sec:ver}

All code is in \texttt{proofs/code/2026-09-16-c2-Q150-H2/}.  The operator engine
(\texttt{independent\_maya.py}) was written from \cite[Conv.~1.2]{Q149} alone and shares no code
with the Q149 or Q150 engines; it acts on Maya sets directly, so a state need not be given as a
small partition.

\paragraph{V1 (calibration of the engine).}  \cite[Lemma 1.3]{Q149} re-derived against Young
diagrams: $1375$ bead moves ($|\lambda|\le8$, $e\le5$), $0$ mismatches.

\paragraph{V2 (\cite[Thm.~6.7]{Q150c1}, and controls).}  $[R_3^{W},R_6^{\bar W}]$ vanishes on all
$195$ partitions with $|\lambda|\le11$.  Negative controls over the $97$ partitions with
$|\lambda|\le9$, all of which fire: NC1 $W_{(1,2)}{=}1$, $35$ nonzero; NC2 $W_{(3)}{=}1$, $32$;
NC3 $W_{(1,1,1)}{=}1$, $32$; NC4 sign flip on $\bar W_{(2,1,2,1)}$, $2$; NC5 drop
$\bar W_{(2,1,2,1)}$, $2$; NC6 drop $\bar W_{(2,3,1)}$, $2$; NC7 add $\bar W_{(6)}{=}1$, $48$.
Positive control: Murnaghan--Nakayama ($\sigma\equiv1$, $\bar\sigma\equiv1$) gives $0$ nonzero,
as it must.  NC1 is the control that located the omission of \S\ref{sec:verify36}.

\paragraph{V3 (admissible supports).}  All $Z\subseteq\cube^{n}$ satisfying \eqref{eq:S} were
enumerated for $d\le7$ by constraint propagation, and independently by brute force over all
$2^{2^{n}}$ subsets for $d\le5$.  The two agree: $4,8,26$ admissible sets at $d=3,4,5$.  (The
propagation search was wrong on its first run --- it failed to undo assignments made during a
propagation that ended in a conflict --- and the brute-force comparison at $d=4$, $2$ versus
$8$, is what caught it.)  Totals, and proper nonempty counts:
\[
\begin{array}{r|ccccc}
 d & 3 & 4 & 5 & 6 & 7\\\hline
 \text{admissible} & 4 & 8 & 26 & 66 & 722\\
 \text{proper}\ne\emptyset & 2 & 6 & 24 & 64 & 720
\end{array}
\]
For each of these, Theorem~\ref{thm:struct} was checked directly: every proper nonempty $Z$ lies
in some $\Cset_{j,b}$ ($720/720$ at $d=7$), and for the minimal such $j_0$ the asserted
flip-invariances hold in every case.  The distribution of $j_0$ at $d=7$ is $712$ at $j_0=1$,
$6$ at $j_0=2$, $2$ at $j_0=3$; the $j_0$ values factor as: $2\cdot2^{d-3}$ sets of size
$2^{d-3}$ are the cylinders themselves, one for each of the $2\lfloor(d-1)/2\rfloor$ distinct
cylinders, and at $d=7$ that is $6$ sets of size $16$ --- matching the $6$ distinct
$\Cset_{j,b}$.  Also checked: $\one_Z$ solves \eqref{eq:2Bstar} for every admissible $Z$
(Proposition~\ref{prop:supp}), all $d\le7$.

\paragraph{V4 (Theorems~\ref{thm:singleton} and \ref{thm:singleton2}).}  For every $d\le9$ and
every $w\in\cube^{d-1}$ the criteria \eqref{eq:NB} and ``antipalindrome'' were compared against
a direct evaluation of \eqref{eq:2Bstar}, resp.\ \eqref{eq:2Bss}, over all instances.  Exact
agreement, both conditions, all $d$.  Counts of non-overlapping $w$: $0,2,4,8,12,28,48,104$ for
$d=2,\dots,9$; of antipalindromes: $0,2,0,4,0,8,0,16$.  Note the counts are $2^{d-2}$ (the
$p=0$ clause alone) exactly for $d\le5$ and strictly smaller from $d=6$ on, which is the content
of Remark~\ref{rem:endrows}.

\paragraph{V5 (Theorem~\ref{thm:main} end to end, with a liveness check).}  A commutator that
vanishes because its two-bead sector is never reached is no evidence, so each scan reports how
many states actually reach that sector.  Exhaustive over all $2^{13}$ (resp.\ $2^{15}$) Maya
windows: at $(e,f)=(3,6)$, $w\in\{01,10\}$, $0$ nonzero commutators with $187$ resp.\ $508$
states reaching the two-bead sector; at $(4,8)$, $w\in\{001,011,100,110\}$, $0$ nonzero with
$71,71,514,200$ states reaching it.  Corollary~\ref{cor:48} is the case $w=001$.  Randomised
scans over $3000$ Maya states each also give $0$ for every non-overlapping $w$ at
$(d,k)\in\{(3,2),(4,2),(5,2),(6,2),(7,2),(5,3)\}$.

\paragraph{V5$'$ (the negative control for Theorem~\ref{thm:main}, constructed not sampled).}
Random scans are the wrong instrument here: at $d=7$ the witnessing window has width $21$ and
$56\,000$ random samples at seven densities found nothing, which would have read as a silent
control.  Instead the witness is built from the geometry of the two-bead element \cite[Lemma 4.3]{Q149}: bead $A$ at site $0$
moving $0\mapsto d$, bead $B$ at site $\delta$ moving $\delta\mapsto\delta+kd$, with the window
filled so that the $\sigma$-word is $w$ and the $\bar\sigma$-word is $w\,0\,w\cdots0\,w$.  For
\emph{every} $w$ with $w_1\ne w_{d-1}$ that fails \eqref{eq:NB} --- the hard cases, the ones the
$p=0$ clause does not catch --- and for every $d\le10$, the constructed state has a nonzero
two-bead commutator element: $112/112$ fire.  The same construction applied to
non-overlapping $w$ gives zero at every $\delta$, for all $d\le7$.  (An earlier version of this
check used $\delta=n+1-p$ instead of $\delta=n-p$ and reported zero everywhere; the fact that
the \emph{positive} cases also reported zero is what exposed it.)

\emph{This control is claim-guided, and I flag it as such.}  The construction chooses $\delta$
using the failure condition of Theorem~\ref{thm:singleton} itself, so it is not independent of
the statement it tests: what it confirms is that the proof's bookkeeping survives contact with
the actual operator, not that the criterion is right as seen from outside.  (My trajectory
validator marks it circular, correctly.)  The independent counterpart is a blind random scan of
Maya states, which knows nothing of the criterion: at $d=6$ each of the four words with
$w_1\ne w_{d-1}$ that fail \eqref{eq:NB} produces $1$--$4$ nonzero commutators in $3000$
samples, while every non-overlapping $w$ produces $0$.  That check does not reach $d\ge7$,
because the witnessing window grows faster than a random sample will find it --- which is why
the constructed control exists at all.

\paragraph{V7 (generation).}  Theorem~\ref{thm:closure} was checked against a direct
evaluation of \eqref{eq:S}: over \emph{all} $2^{2^{n}}$ subsets for $d\le4$, and over the
admissible sets together with all $2$-element sets for $d=5,6$.  No mismatches.  Intersection
closure (Proposition~\ref{prop:moore}) was checked on every pair of admissible sets for $d\le6$.
The closure of every single word was computed for $d\le6$; sizes as tabulated in
\S\ref{sec:closure}.

\paragraph{V6 (values).}  For every admissible $Z$ at $d\le6$ the lattice $L_Z$ of
\S\ref{sec:values} was put in Smith normal form.  Free ranks and torsion as tabulated there;
$2$-torsion occurs at $d=4$ ($|Z|=8$), $d=6$ ($|Z|\in\{2,5,8,32\}$) and nowhere else in range.

\paragraph{Scope, against the quantifier.}  Theorems~\ref{thm:singleton},
\ref{thm:singleton2}, \ref{thm:struct} and \ref{thm:main} are proved for all $d$; the
computations above are confirmations, not the proof.  The quantifier ranges checked are: all
$w\in\cube^{d-1}$ for $d\le9$ (V4); all $Z\subseteq\cube^{d-1}$ for $d\le7$ (V3); all Maya
windows of width $13$, $15$ for the two smallest $(e,f)$ (V5); all hard failing $w$ for
$d\le10$ (V5$'$).  Nothing is checked for $k\ge4$, and nothing for $e>d$ --- the latter is
exactly gap (H1) of \cite{Q150c1}, untouched here.

\section{What is and is not proved}\label{sec:gaps}

\subsection*{Proved}
\begin{enumerate}
\item Proposition~\ref{prop:supp}: admissible supports $=$ supports of solutions, with $\one_Z$
  as base point.
\item Theorem~\ref{thm:singleton}: $c\,\one_w$ satisfies (2B$^*_d$) iff $w$ is non-overlapping.
  Unconditional, one domain appeal, witness $c$.
\item Theorem~\ref{thm:singleton2}: $c\,\one_w$ satisfies (2B$^{**}_d$) iff $w$ is an
  antipalindrome; none for even $d$.
\item Theorem~\ref{thm:main} and Corollary~\ref{cor:48}: for every $d\ge3$, every
  non-overlapping $w$ and every $k\ge2$, an operator adding the single ribbon shape
  $\alpha(w)$ commutes with an explicit rank-$kd$ operator supported on $2^{k-1}$ shapes.
\item Theorem~\ref{thm:struct}: the containment half of (H2), with the flip-invariance
  refinement; and Corollary~\ref{cor:three}, which recovers the $\gcd\le2$ nonvanishing of
  \cite[Thm.~6.1]{Q150c1} at the level of supports and explains the threshold.
\item Corollaries~\ref{cor:maximal} and \ref{cor:quarter}: the cylinders are exactly the maximal
  proper supports, whence the quarter dichotomy --- a nonzero commuting weight is nowhere zero
  or supported on at most a quarter (rank $d$), resp.\ $4^{-e/d}$ (rank $e$), of the shapes.
\item Theorem~\ref{thm:closure} and Proposition~\ref{prop:moore}: $Z$ is admissible iff closed
  under the linked-pair flip rule \eqref{eq:rule}; the admissible sets form a closure system,
  and Theorem~\ref{thm:singleton} is the statement $\langle\{w\}\rangle=\{w\}$.
\item \S\ref{sec:verify36}: independent re-verification of \cite[Thm.~6.7]{Q150c1} and the
  correction $W_{(1,2)}=0$ to its statement.
\end{enumerate}

\subsection*{Not proved}
\begin{enumerate}
\item[(H2a)] \textbf{Which subsets of a cylinder are admissible.}  Theorem~\ref{thm:struct}
  reduces the classification to: for each $j$ with $2j\ne d$, describe the admissible
  $Z\subseteq\Cset_{j,b}$ that are flip-invariant at $\{1,\dots,j-1\}\cup\{d-j+1,\dots,d-1\}$.
  Such a $Z$ is the same thing as a set $S\subseteq\cube^{\,d-2j-1}$ of \emph{middle} words, and
  unwinding \eqref{eq:S} for $\delta$ in the remaining range turns the condition into: for
  $1\le i\le m-1$ ($m=d-2j-1$) and all $p,r\in\cube^{\,i-1}$, $q\in\cube^{\,m-i-1}$,
  \[
     \bigl[(p,c,b,q)\in S\ \wedge\ (q,1-b,c,r)\in S\bigr]\quad
     \text{is the same for }c=0\text{ and }c=1 .
  \]
  This is \eqref{eq:S} again but with a pinned letter adjacent to the toggled one on each side
  and a shared block one shorter --- structurally similar, not literally the same, so it does
  not self-reduce.  I verified the reduction at $j=1$, $d\le6$ but do not have a closed
  description of its solutions.  The counts of admissible $S$ are $1,3,11,31,356$ for
  $m=0,\dots,4$, which I do not recognise; the jump $31\to356$ suggests the answer is not a
  product formula.  \emph{Theorem~\ref{thm:singleton} is the case $|S|=1$ of this question,
  solved.}  By Theorem~\ref{thm:closure} and Proposition~\ref{prop:moore} the question has been
  turned into a concrete one about an explicit closure operator --- ``describe the closed sets
  of \eqref{eq:rule}'' --- rather than about a system of equations; that is progress of a kind,
  but it is not the answer.

  \textbf{And the obvious induction does not continue.}  Theorem~\ref{thm:struct} works because
  the maximal proper admissible sets have a uniform shape (Corollary~\ref{cor:maximal}: they are
  the cylinders, all of size $2^{d-3}$).  One level down this fails.  Computing the maximal
  proper admissible subsets \emph{of} $\Cset_{1,0}$ in middle-word coordinates gives
  \[
  \begin{array}{r|l|l}
    m=d-3 & \#\text{ maximal} & \text{their sizes (out of } 2^{m})\\\hline
    2 & 3  & 2,3,3\\
    3 & 4  & 4,4,5,5\\
    4 & 11 & 4,4,5,5,5,5,6,6,7,9,13
  \end{array}
  \]
  None is a sub-cylinder $\{x_i=a,\ x_{m+1-i}=1-a\}$, and at $m\ge3$ none is a co-singleton.
  There is no uniform maximal shape to induct on, so the argument of
  Theorem~\ref{thm:struct} has no second step.  I record this as a \emph{negative} result: it
  is the route a next attempt would try first, and it is closed.
\item[(H2b)] \textbf{A uniform description of the value lattice $L_Z$.}  \S\ref{sec:values}
  computes it case by case; I have no formula for its rank in terms of $Z$, and the sporadic
  $2$-torsion is unexplained.
\item[(H1)] \textbf{Untouched.}  The middle condition of \cite[Rem.~5.2]{Q150c1} --- which
  $\tau$ have $\tau^{\otimes k}$ satisfying (2B$^*_{kd}$) --- is not addressed.  Note however
  that Theorem~\ref{thm:singleton2} gives the first clean statement about it from the (2B$^{**}$)
  side: for even $d$ \emph{no} single-shape $\tau$ satisfies (2B$^{**}_d$), yet by
  Theorem~\ref{thm:main} single-shape $\tau$ at $d=4,6,8$ do produce commuting pairs at
  $e=d$.  So for even $d$ the route through (2B$^{**}$) of \cite[Thm.~6.6]{Q150c1} misses the
  single-shape examples entirely --- which is why they are new.
\end{enumerate}

\subsection*{Effect on \cite{Q150c1}}

No theorem of \cite{Q150c1} changes.  Theorem~\ref{thm:struct} upgrades its computational
observation (``the admissible supports at $d\le5$ are the full support and the subsets of the
$\Cset_{\delta,b}$'', \S8 (H2)) to a theorem for all $d$, and Corollary~\ref{cor:three} gives an
independent proof of the support half of its Theorem~6.1.  Theorem~\ref{thm:main} enlarges the
family of counterexamples in its Theorem~6.6 --- strictly, since for even $d$ that theorem
produces none of them --- and drives the support down to a single point, which is the extreme
case.  The only textual change required is the omitted value in its Theorem~6.7,
\S\ref{sec:verify36}.

\begin{thebibliography}{9}
\bibitem{Q149} Clio, \emph{Commuting ribbon operators with shape-dependent weights},
  2026-09-15, \texttt{proofs/2026-09-15-c1-Q149-shape-weights.tex}.
\bibitem{Q150c1} Clio, \emph{The nonvanishing lemma for ribbon shape weights is false, and the
  true statement is a commuting-words theorem}, 2026-09-16 (cycle 1),
  \texttt{proofs/2026-09-16-c1-Q150-nonvanishing.tex}.
\end{thebibliography}

\end{document}

